\chapter{Introduzione}

Il presente capitolo fornisce una panoramica del progetto di tirocinio, offrendo una spiegazione sull'ambito del progetto e indicando le tecnologie impiegate per la sua realizzazione.

\section{Introduzione al progetto}

Il progetto di tirocinio mira allo sviluppo e l'integrazione di un prototipo di applicazione per la visualizzazione di mappe geospaziali,
avvalendosi degli standard dell'\textit{Open Geospatial Consortium} (\textit{OGC}) e delle più moderne tecnologie web.
\\Tale iniziativa si inserisce nel contesto dell'importanza sempre crescente dei dati
geospaziali nelle applicazioni web e della necessità di fornire strumenti interattivi e intuitivi per la loro visualizzazione e manipolazione.
\\Il lavoro si concentra sia sull'implementazione di un'applicazione frontend per la visualizzazione di mappe geospaziali, includendo un'interfaccia utente
interattiva, sia sullo sviluppo di un backend per la gestione e la fornitura dei dati delle mappe.
\vspace{\baselineskip}Lo scopo principale del progetto di tirocinio è stato quello di contribuire allo sviluppo di una applicazione chiamata Inspicio attraverso l'implementazione di una funzionalità di visualizzazione delle mappe geospaziali.
Questo aggiornamento ha permesso agli utenti di esplorare interattivamente i ponti in Italia e di accedere a informazioni cruciali come eventi sismici, frane e altri dati correlati al territorio. L'integrazione delle mappe ha arricchito notevolmente le capacità di analisi e di monitoraggio dell'applicazione, migliorando così la sua utilità e la sua rilevanza nel contesto delle infrastrutture nazionali.